

\documentclass[12pt,oneside]{article}

%%%%%%%%%%%%%%%%%%%%%%%%%%%%
%%   Zusaetzliche Pakete  %%
%%%%%%%%%%%%%%%%%%%%%%%%%%%%
\usepackage{enumerate}  
\usepackage{fancyhdr}
\usepackage{a4wide}
\usepackage[utf8]{inputenc}
\usepackage{graphicx}
\usepackage{palatino}
\usepackage{multirow}
\usepackage{booktabs}
\usepackage{titlesec}
\usepackage{amsthm}
\usepackage{amssymb}
\usepackage{amsmath}
\usepackage[utf8]{inputenc}
\usepackage[options ]{algorithm2e}
\usepackage[ruled,vlined]{algorithm2e}
\usepackage{lmodern}
\usepackage[utf8]{inputenc}
\usepackage[english,german]{babel}
\usepackage{xcolor}
\usepackage{amsfonts}
\usepackage{etoolbox}
\usepackage[]{algorithm2e}
\usepackage{blindtext}
\usepackage{algorithm}
\usepackage[noend]{algpseudocode}
\usepackage{enumitem}% http://ctan.org/pkg/enumitem

% notendig für Definitionen und Theoreme
\newtheorem{theorem}{Theorem}[section]
\newtheorem{corollary}{Corollary}[theorem]
\newtheorem{lemma}[theorem]{Lemma}

\newtheorem{prop}{Proposition}[section]

\theoremstyle{remark}
\newtheorem*{remark}{Remark}

\theoremstyle{definition}
\newtheorem{definition}{Definition}[section]
%%%%%%%%%%%%%%%%%%%%%%%%%%%%%%

%folgende Zeile auskommentieren für englische Arbeiten
\usepackage[ngerman]{babel}
%folgende Zeile auskommentieren für deutsche Arbeiten
%\usepackage[ngerman, english]{babel}

\usepackage[T1]{fontenc}
\usepackage[utf8]{inputenc}
\usepackage[bookmarks]{hyperref}
\usepackage[justification=centering]{caption}
\usepackage[style=authoryear,natbib=true,backend=biber,maxbibnames=20]{biblatex}
\usepackage{csquotes}
\bibliography{literatur}

\setlength{\parindent}{0em} 
\setlist[itemize]{noitemsep, topsep=0pt}
\setlist[enumerate]{noitemsep, topsep=0pt}


%%%%%%%%%%%%%%%%%%%%%%%%%%%%%%
%% Definition der Kopfzeile %%
%%%%%%%%%%%%%%%%%%%%%%%%%%%%%%

\pagestyle{fancy}
\fancyhf{}
\cfoot{\thepage}
\setlength{\headheight}{16pt}

%%%%%%%%%%%%%%%%%%%%%%%%%%%%%%%%%%%%%%%%%%%%%%%%%%%%%
%%  Definition des Deckblattes und der Titelseite  %%
%%%%%%%%%%%%%%%%%%%%%%%%%%%%%%%%%%%%%%%%%%%%%%%%%%%%%

\newcommand{\JMUTitle}[9]{

  \thispagestyle{empty}
  \vspace*{\stretch{1}}
  {\parindent0cm
  \rule{\linewidth}{.7ex}}
  \begin{flushright}
    \vspace*{\stretch{1}}
    \sffamily\bfseries\Huge
    #1\\
    \vspace*{\stretch{1}}
    \sffamily\bfseries\large
    #2\\
    \vspace*{\stretch{1}}
    \sffamily\bfseries\small
    #3
  \end{flushright}
  \rule{\linewidth}{.7ex}

  \vspace*{\stretch{1}}
  \begin{center}
    \includegraphics[width=4in]{logo} \\
    \vspace*{\stretch{1}}
    \Large  Bachelorarbeit   \\

    \vspace*{\stretch{2}}
   \large Lehrstuhl für Mathematik \\
    \large und Informatik \\
    \large Universität Leipzig\\
    \vspace*{\stretch{1}}
    \large Betreuer:  #8 \\[1mm]
    
    \vspace*{\stretch{1}}
    \large Leipzig, den #7
  \end{center}
}

\titlespacing*{\section}
{0pt}{3.5ex plus 1ex minus .2ex}{.2ex plus .2ex}
\titlespacing*{\subsection}
{0pt}{1.5ex plus 1ex minus .2ex}{.2ex plus .2ex}
\titlespacing*{\subsubsection}
{0pt}{1.5ex plus 1ex minus .2ex}{.2ex plus .2ex}

%%%%%%%%%%%%%%%%%%%%%%%%%%%%
%%  Beginn des Dokuments  %%
%%%%%%%%%%%%%%%%%%%%%%%%%%%%

\begin{document}

  \JMUTitle
      {Komplexität und Evaluation des AKS-Primzahltests  }  
      {Salman Salman}                        
      {3753924}
      
      {Fakultät für Informatik und Mathematik}  % Name der Fakultaet
      {Leipzig 2020}                          % Ort und Jahr der Erstellung
      {\today}                              % Tag der Abgabe
      {Prof. Dr. Andreas Maletti}               % Name des Erstgutachters
      {Zweitgutachter}                          % Name des Zweitgutachters
      
  \clearpage

\lhead{}
\pagenumbering{Roman} 
    \setcounter{page}{1}

\tableofcontents
\clearpage

\addcontentsline{toc}{section}{\listfigurename}
\listoffigures

\addcontentsline{toc}{section}{\listtablename}
\listoftables
\clearpage

\setlength{\parskip}{0.5em} 


%%%%%%%%%%%%%%%%%%%%%%%%%%%%
%%  Kurzzusammenfassung   %%
%%%%%%%%%%%%%%%%%%%%%%%%%%%%
\lhead{Abstract}
\section*{Abstract}
Primzahlen haben in der Informatik, speziell im Anwendungsgebiet der Kryptographie, eine sehr hohe Relevanz für moderne kryptographische Systeme ist es von Wichtigkeit selbige schnell bestimmen zu können. Hierzu werden effiziente Algorithmen zur Lösung des mathematischen Primalitätsproblem benötigt. Das Primalitätsproblem umfasst die Frage um die Entscheidung, ob eine gegebene Zahl eine Primzahl ist oder nicht. Der erste deterministische Primzahltest in Polynomialzeit wurde von den indischen Informatikern Agrawal, Kayal, und Saxena vorgestellt. Der nach ihnen benannte AKS-Algorithmus wird im Rahmen dieser Arbeit repräsentiert, implementiert und evaluiert. Außerdem wird die Korrektheit des Algorithmus experimentell geprüft und mathematisch bewiesen. 


%%%%%%%%%%%%%%%%%%%%%%%%%%%%
%%  Einstellungen  %%
%%%%%%%%%%%%%%%%%%%%%%%%%%%%
\clearpage
\pagenumbering{arabic}  
    \setcounter{page}{1}
\lhead{\nouppercase{\leftmark}}

%%%%%%%%%%%%%%%%%%%%%%%%%%%%
%%  Hauptteil  %%
%%%%%%%%%%%%%%%%%%%%%%%%%%%%

 

\section{Einleitung} \label{einleitung}

\subsection{Gliederung der Arbeit}
Die Arbeit gliedert sich in vier Bestandteile, wobei der
erste Teil die Einleitung darstellt. Dort wird ein kurzer Überblick in das Thema gegeben und die Zielsetzung der Arbeit definiert.

Das zweite Kapitel enthält die Grundlagen und die Definitionen. Hier werden Terminologie und Konzepte erläutert, die für das weitere Verständnis der Arbeit eine Rolle spielen. 

Der Algorithmus und seine Korrektheit sind die Bestandteile des dritten Kapitels. Zuerst wird die Grundidee des Algorithmus erläutert, danach wird der Algorithmus in seiner ursprünglichen Form(wie im AKS-Paper) angegeben. Anschließend wird der Korrektheitsbeweis des Algorithmus geführt.

Im vierten Kapitel wird die Laufzeitanalyse des Algorithmus vorgestellt, dabei werden die Laufzeit der einzelnen Schritte (STEPS 1-6) des Algorithmus ausführlich analysiert.    

\subsection{Überblick}
Eine natürliche Zahl $n \geq 2$ ist eine Primzahl, wenn keine der Zahlen $r \in (1,n)$ $n$ teilt, also $\forall r \in (1,n), r \nmid n$. Die Frage, ob eine Zahl eine Primzahl ist oder nicht, war schon im antiken Griechenland interessant. Euklid hat sich mit dieser Frage beschäftigt und beweisen, dass es unendlich viele Primzahlen gibt. Im 3. Jahrhundert v. Chr. hat der griechische Mathematiker Eratosthenes einen Algorithmus zur Bestimmung aller Primzahlen von 1 bis zu einer gegebenen Zahl $n $ vorgestellt, dieser Algorithmus ist jedoch zur Lösung des Primalitätsproblem nicht effizient, da wir uns nur dafür interessieren, ob die Zahl $n$ prim ist oder nicht. Im 17. und 18. Jahrhundert haben Wissenschaftler versucht eine Formel für Primzahlen zu erfinden. Der französische Wissenschaftler Marin Mersenne hat eine Formel für die sogenannten Mersenne-Zahlen entwickelt, 51 der Mersenne-Zahlen sind Primzahlen. Solche Zahlen heißen Mersenne-Primzahlen, die größte Primzahl heute $2^{82,589,933} − 1$ ist eine Mersenne-Primzahl. Der Algorithmus von Eratosthenes braucht länger als das Alter der Sonne, um zu entscheiden ob die obere Mersenne-Zahl eine Primzahl ist oder nicht. Deswegen ist es sinnvoll effizientere Algorithmen zu entwickeln. Ein anderer französischer Mathematiker Pierre de Fermat hat eine Gleichung zur Lösung des Primalitätsproblems gefunden, die für alle Primzahlen wahr ist, auf seine Gleichung bauen viele Primzahltests auf. \newline\newline
Im Laufe der Zeit wurden auch mehrere Algorithmen zur Lösung des Primalitätsproblems entwickelt, die entweder von unbewiesenen Hypothesen, wie zum Beispiel die verallgemeinerte Riemannsche Vermutung abhängig waren oder probabilistisch waren. Probabilistische Algorithmen sind randomisierte Algorithmen, die auch ein falsches Ergebnis liefern können. Ein bekannter probabilistischer Algorithmus ist der Miller-Rabin-Primzahltest, er kann das Primalitätsproblem in $O(k log^3 n)$ lösen. Dennoch kann der beste nicht probabilistische Algorithmus (vor AKS) das Primalitätsproblem in $ \Omega(\sqrt{n}) $ Schritten lösen, wobei $n$ die Eingabegröße ist(Anzahl der nötigen Bits, um die Zahl $n$ zu repräsentieren). Solcher Algorithmus braucht exponentielle Zeit, um das Problem zu lösen. In 2002 haben drei Informatiker Agrawal, Kayal, und Saxena den ersten deterministischen unbedingten Algorithmus entwickelt, der das Primalitätsproblem in Polynomialzeit lösen kann. Das heißt sie haben gezeigt, dass das Primalitätsproblem(PRIMES) zur Komplexitätsklasse $P$ gehört.      

\subsection{Zielsetzung}
In dieser Arbeit wird der AKS-Primzahltest Algorithmus aus dem Orginalpaper behandelt. Ziel dieser Arbeit ist die Korrektheit des AKS-Algorithmus zu testen und evaluieren. Dabei wird die Idee Algorithmus dargestellt und erläutert, danach kommt die Analyse der Korrektheit, da wird die Korrektheit zuerst mathematisch bewiesen, dafür wird ein bidirektionaler Beweis geführt, wo gezeigt wird, dass wenn der Algorithmus PRIME liefert, ist die Zahl auch eine Primzahl. Aber auch wenn die Eingabe eine Primzahl ist, liefert der Algorithmus PRIME. Dazu wird seine Korrektheit auf einem Rechner implementiert und getestet. Der Algorithmus ist dafür bekannt, das Primalitätsproblem in Polynomialzeit zu lösen, diese Eigenschaft wird auch zuerst mathematisch bewiesen. Schließlich soll die Polynomialzeit Eigenschaft durch eine Grafik experimentell bewiesen.  

Um die Korrektheit zu demonstrieren, soll der Algorithmus in Python implementiert werden, dabei wird die Ausgabe des Algorithmus bei mehreren Zahlen(meistens große Zahlen) geprüft. Zudem soll auch die Laufzeit durch eine Grafik illustriert werden.  


\section{Definitionen und Grundlagen}
In diesem Kapitel werden die essenziellen Begriffe und Theoreme aus der Zahlentheorie, der Algebra und der Theorie der zkyklotomischen Polynome, die für den AKS-Algorithmus relevant sind, definiert. In diesem Kapitel handelt es sich darum, einen formalen Literaturhinweis zu haben, um später die hier definierten Begriffe und Theoreme zu referenzieren.
% Definitionen der Zhalentheorie 
% Definitonen
\subsection{Zahlentheorie}
\theoremstyle{definition}
\begin{definition}\label{Df_1}
Seien $a,b \in \mathbb{N}$. Der \textbf{größte gemeinsame Teiler} von $a$ und $b$ wird mit $(a,b)$ bezeichnet, ist die größte positive Zahl $n$, sodass $n \mid a$ und $n$ $ \mid b$
\end{definition}

\smallskip 

\begin{definition}\label{Df_2}
Zwei Zahlen $a,b \in \mathbb{N}$ heißen genau dann \textbf{Teilerfremd}, wenn $(a,b) = 1$.
\end{definition}

\smallskip

\begin{definition}
Seien $n, m$ zwei natürliche Zahlen, dann heißt die kleinste positive natürliche Zahl, die sowohl ein Vielfaches von $n$, als auch von $m$ \textbf{kleinstes gemeinsames Vielfaches} beider zahlen. Hier wird das kleinste gemeinsame Vielfaches mit \textbf{LCM} bezeichnet.
\end{definition}

\smallskip 

\begin{definition}\label{Df_3}
Seien $a, b, n \in \mathbb{N}$. a ist genau dann \textbf{kongruent} zu $b$ modulo $n$, wenn $n \mid a - b $ gilt, dies wird mit $a = b$ (mod $n$) bezeichnet.  
\end{definition}

\smallskip 

\begin{definition}\label{Df_4}
Seien $r,n \in \mathbb{N}$ mit $(n,r) = 1$, dann ist die \textbf{Ordnung} von n modulo r das kleinste k, sodass $n^k = 1 (mod $ r). Die Ordnung wird mit $o_{r}(n)$ bezeichnet.
\end{definition}

\smallskip


\begin{definition}\label{Df_5}
Sei $n \in \mathbb{N}$, die \textbf{Primzahlzerlegung} von $n$ ist die Darstellung der Zahl als Produkt ihrer Primfaktoren \newline
\begin{equation}
    n = p_{1}^{e_{1}}p_{2}^{e_{2}}...p_{M}^{e_{M}}
\end{equation}
Wobei $e_{k}$ die Vielfachheit der Primzahl $p_{k}$ ist.
\end{definition}

\smallskip


\begin{definition}\label{Df_6}
Sei $n \in \mathbb{N}$ mit $n > 1$. Die \textbf{eulersche Phi-Funktion} wird mit $\phi(n)$ bezeichnet, ist die Anzahl der Zahlen $k, \, 1 \leq k \leq n$, sodass $(k,n) = 1$
\end{definition}

\smallskip

\begin{theorem}[\textbf{Gaußsche Summeformel}]\label{gauss}
Sei n eine natürlich Zahl, dann gilt:
\begin{equation}
    \sum_{i = 1}^{n} i = \frac{n (n + 1)}{2} = \frac{n^2 + n}{2} 
\end{equation}
\end{theorem}

\smallskip
% Theoreme
\begin{theorem}[\textbf{Satz von Euler}]\label{Th_1}
Seien $a,n \in \mathbb{N}$ teilerfremd, \newline dann gilt $a^{\phi(n)} = 1 $(mod n).
\end{theorem}

\begin{proof}
    
Seien $R =\{ a \in \mathbb{N} \mid (a,n) = 1 \} =  \{a_{1}, a_{2},...,a_{\phi(n)} \}$ und $c \in R, c \leq n$.\newline\newline Da $(c,n) = (a_{i},n) = 1$, gilt auch $(c \cdot a_{i}, n) = 1$ \newline\newline und damit ist auch $a_{i} = c \cdot a_{j}$(mod $n$) wahr für ein $a_{j} \in R \Rightarrow a_{i} = a_{j}$. Das heißt es gilt:\newline
\begin{equation}
     \prod_{k = 0}^{\phi(n)} c \cdot a_{k} =
     c^{\phi(n)}\prod_{k = 0}^{\phi(n)} a_{k} = \prod_{k = 0}^{\phi(n)}  a_{k} (mod \, n )
\end{equation}
    
     $\Rightarrow c^{\phi(n)} = 1$ (mod $n$).

\end{proof}

\smallskip

\begin{theorem}[\textbf{Kleiner fermatscher Satz}]\label{Th_2}
Sei $a \in \mathbb{N}$ und p eine Primzahl dann gilt:\newline
\begin{enumerate}
    \item a und p sind genau dann Teilerfremd, wenn $a^{p-1} = 1 $ (mod p) gilt.
    \newline
    \item $\forall a $ gilt  $a^p = a$ (mod p)
\end{enumerate}
\end{theorem}

%Theoretische Informatik 

\subsection{Komplexitätstheorie}

Hier handelt es sich darum, die relevanten Begriffe und aus dem Gebiet der Komplexitätstheorie, eindeutig zu definieren. Grundlagen der Komplexitätstheorie, wie Turingmaschinen, Polynomiallaufzeit, und Entscheidungsprobleme werden in dieser Arbeit nicht definiert, jedoch sind sie eine Voraussetzung, um die hier behandelten Themen zu verstehen.   

\begin{definition}
In der Komplexitätstheorie ist \textbf{P} die Komplexitätsklasse, die alle Entscheidungsprobleme enthält, die in Polynomialzeit für deterministische Turingmaschinen lösbar sind. 
\end{definition}

\begin{definition}
Eine Funktion $f(n)$ ist $ O^{\sim}(t(n))$, wenn
\begin{equation}
    f(n) = O(t(n) \cdot poly(log t(n)))
\end{equation}
\textbf{Beobachtung:}
\begin{equation}
    O^{\sim}(log^k n) = O(log^k n) \cdot poly(log log^k n) = O(log^k n) \cdot poly(log( k \cdot log n)) = O(log^{k+\epsilon}n)
\end{equation}
Für alle $\epsilon > 0$, wenn der Algorithmus eine Laufzeit von $O^{\sim}(log^k n)$ hat. Dann läuft er in Polynomialzeit von $log n$.
\end{definition}

% Algerba

\subsection{Algebra}
Es wird davon ausgegangen, dass die grundlegenden Definitionen von Gruppen, Ringen, Körpern und ihre Eigenschaften dem Leser bekannt sind. Hier werden Theoreme und Begriffe der Algebra und Polynomentheorie, die später für den Korrektheitsbeweis von Bedeutung sind, vorgestellt und bewiesen.\newline

\begin{definition}
Sei $R$ ein Ring, dann ist der \textbf{Polynomring} $R[X]$ die Menge aller Polynome der Form $a_{0} + a_{1}X + a_{2} X^2 + ... + a_{n}X^n$, wobei $a_{0},a_{1}...,a_{n} \in R$.
\end{definition}

\smallskip

\begin{definition}
Sei $I$ eine Teilmenge des Rings $(R,+,\cdot)$. $I$ heißt \textbf{Ideal}, wenn gilt:
\begin{enumerate}
    \item $0 \in I$.
    \item $\forall a,b \in I$ ist $a + b \in I$.
    \item $\forall a \in I$ und $r \in R$ ist $r \cdot a \in I$.
    \item $\forall a \in I$ und $r \in R$ ist $a \cdot r \in I$.
\end{enumerate}
\end{definition}

\smallskip 

\begin{definition}
ist $(R, +, \cdot)$ ein Ring und $I$ ein Ideal von R, dann bildet die Menge $R/I = \{ a + I \mid a \in R\}$ der Äquivalenzklassen modulo $I$ mit folgenden Verknüpfungen einen Ring:
\begin{itemize}
    \item $(a + I) + (b + I) = (a + b) + I$.
    \item $(a + I) \cdot (b + I) = (a \cdot b ) + I$.\newline
\end{itemize}
Diesen Ring nennt man den \textbf{Quotientenring}. 
\end{definition}

\smallskip

\begin{theorem}\label{th_25}
Sei n eine Primzahl, dann gilt ${n \choose i} = 0 $(mod n).
\end{theorem}

\begin{proof}
\begin{equation}\label{modb}
    {n \choose i} = \frac{(n - i - 1) \cdot \cdot \cdot (n - 1) \cdot n }{i!}
\end{equation}
Da $n$ im Zähler steht und n eine Primzahl ist(nicht durch Zahlen im Nenner teilbar), muss die obere Gleichung (\ref{modb}) durch n teilbar sein. 
\end{proof}

\smallskip
\begin{flushleft}
\begin{theorem}
${2n + 1 \choose n} > 2^{n+1}, \forall n \geq 2$
\end{theorem}

\begin{proof}
Induktion: Sei n = 2, offensichtlich gilt ${5 \choose 2} > 2^3$. Sei nun das obere Theorem für ein beliebiges $k > 2 $ wahr.\newline\newline
für $n = k + 1$ ist das folgende zu zeigen: 

\begin{equation}\label{ind_1}
    \frac{(2k + 3)!}{(k + 2)!\cdot(k + 1)!} > 2^{k+2}.
\end{equation}
\newline\newline
Die obere Ungleichung (\ref{ind_1}) lässt sich wie folgt umschreiben:\newline\newline


\begin{equation}
     \underbrace{\frac{(2k + 1)!}{(k + 1)! \cdot k!}}_{nach IH > 2^{k + 1}} \cdot \frac{(2k + 2) \cdot (2k + 3)}{(k + 2) \cdot (k + 1)} > 2^{k+1} \cdot 2.\newline\newline
\end{equation}

Der erste Teil der . Multiplikation auf der linken ist nach der Induktionshypothese größer als $2^{k+1}$. Es bleibt nur zu zeigen, dass der rechte Teil der Multiplikation größer als 2 ist, das heißt: \newline\newline

\begin{equation}\label{ind_eq}
\frac{(2k + 2) \cdot (2k + 3)}{(k + 2) \cdot (k + 1)} = 2 \cdot \underbrace{\frac{(2k + 3 )}{k + 2}}_{ > 1} > 2
\end{equation}
Somit ist die Ungleichung für $n = k + 1$ wahr.
\end{proof}

\end{flushleft}
\smallskip 

\begin{theorem}[\textbf{Binomischer Lehrsatz}]\label{Th_3}
Seien R ein kommutativer Ring und n eine natürliche Zahl, dann gilt für $a,b \in $R:\newline\newline
 \begin{equation}
     (a + b)^n  = \sum_{k=0}^n {n \choose k} a^k b^{n-k}
 \end{equation}
\end{theorem}

\smallskip

\begin{theorem}\label{funny_id}
Seien $a,b,n \in \mathbb{N}$ Dann gilt:\newline\newline \begin{equation}
    a^{n} - b^n = (a - b) \cdot \sum_{k = 0}^{n - 1} a^{n -1 - k} b^k.
\end{equation}
\end{theorem}


\smallskip

% Polynome 
\subsection{Zyklotomische Polynome}

\begin{definition}
Sei n eine natürliche Zahl, die n-te \textbf{Einheitswurzel} wird mit $\zeta$ bezeichnet ist eine komplexe Zahl, sodass
\begin{equation}
    \zeta^n = 1.
\end{equation}
Z.B. 1 und -1 sind die quadratischen Einheitswurzeln, und 1, -1, i, -i, sind die Einheitswurzeln für $n = 4$.    
\end{definition}

\smallskip

\begin{definition}
Für jede natürliche Zahl $n$, gibt es $n$ viele Einheitswurzeln, die durch:\newline
\begin{equation}
    e^{2k\pi i/n} = cos(\frac{2k \pi }{n}) + i sin(\frac{2k \pi}{n})
\end{equation}
definiert sind, mit $1 \leq k \leq n$.
\end{definition}

\smallskip

\begin{definition}\label{ord_def}
Sei $\mathbb{K}$ ein Körper, und $a \in \mathbb{K}^x$. Dann ist \textbf{Ordnung} $ ord_{ \mathbb{K} }(a)$ , die kleinste natürliche Zahl $k$, für die $a^k = 1 $ gilt, wenn ein solches $k$ nicht existiert, dann hat a eine unendliche Ordnung. 
\end{definition}

\smallskip

\begin{theorem}
Sei $\mathbb{K}$ ein Körper mit $m$ Elementen, dann gilt:
\begin{equation}
    a^{m-1} = 1
\end{equation}
für alle $a \in \mathbb{K}^x$\newline
%TODO : cite the ref. 
Beweis: siehe (Fields and Cyclotomic Polynomials) 
\end{theorem}

\smallskip

\begin{theorem}
Sei $\mathbb{K}$ ein Körper und $a \in K^x$. Weiterhin sei $ n \in \mathbb{N}$, mit $n \geq 1$. Dann ist $a^n = 1 \Leftrightarrow	 $ $ord_{\mathbb{K}}(a) \mid n$. 
\end{theorem}

\smallskip

\begin{proof}
$"\Rightarrow"$\newline
Sei $ord_{\mathbb{K}}(a) = k $, wenn $k \mid n$, dann $n = mk$, für ein $ m \geq 1$.
\begin{equation}
    a^n = a^{mk} = (a^k)^m. 
\end{equation}
$a^k$ ist definitionsmäßig(\ref{ord_def}) gleich 1. Daher gilt folgendes:

\begin{equation}
    a^n = a^{mk} = (a^k)^m = (1)^m = 1.
\end{equation}
$"\Leftarrow"$\newline
Sei nun $a^m = 1$, Außerdem seien $i,j \in \mathbb{N}$, sodass:\newline
$im + jk = (m,k)$. Dabei ist k die Ordnung des Körpers(\ref{ord_def}).\newline\newline
Weiterhin gilt: 
\begin{equation}\label{gcd_1}
   a^{(m,k)} = a^{im + jk} = (a^m)^i \cdot (a^k)^j = (1)^i (1)^k = 1.
\end{equation}

Aus (\ref{gcd_1}) und Der Definition(\ref{ord_def}) folgt, dass $(m,k) = k$ und somit $k \mid m$.
\end{proof}

\smallskip

\begin{definition}\label{prim_ein}
Die \textbf{n-te primitive Einheitswurzel} ist jede Einheitswurzel $\zeta$, für die $ord_{\mathbb{C}}(\zeta) = n $ gilt. Die Menge aller primitiven n-ten Einheitswurzeln wird mit \textbf{P(n)} bezeichnet. 
\end{definition}

\smallskip

\begin{definition}
\textbf{Zyklotomisches Polynom:} das n-te zyklotomische Polynom $\Phi_{n}$ ist durch:\newline
\begin{equation}
    \Phi_{n}(x) = \prod_{\zeta \in P(n)} (x -\zeta).
\end{equation}

Dabei ist $P(n)$ die Menge aller primitiven n-ten Einheitswurzeln aus der Definition (\ref{prim_ein}). 
\end{definition}

\smallskip

\begin{theorem}
Für jede natürliche Zahl n . gilt:\newline
\begin{equation}
    x^n - 1 = \prod_{d \mid n} \Phi_{d}(n).
\end{equation}
%TO: cite the reference 
Für den Beweis, siehe (Fields and Cyclotomic Polynomials) 
\end{theorem}

\smallskip


%%%%%%%%%%%%%%%%%%%%%%%%%%%%%%%%%%%%%%%%%%%%%%%%%% AKS ALGORITHMUS %%%%%%%%%%%%%%%%%%%%%%%%%%%%%%%%%%%%%%%%%%%%%%%%%%%%%%%%%
\section{Der AKS-Primzahltest}
\subsection{Grundidee des Algorithmus}
Die Idee des Algorithmus ist basiert auf  Verallgemeinerung des kleineren fermatschen Satz.
\begin{flushleft}
\begin{lemma}\label{hauptlemma}
seien $a,n \in \mathbb{N}$ mit a < n und teilerfremd, dann ist n genau dann eine Primzahl, wenn \newline
\begin{equation}\label{eq:1}
\centerline{ $(X + a)^n = X^n +a $(mod n).}.
\end{equation}\newline
Dabei ist X ein Polynom über dem Ring $\mathbb{Z}_{n}[X]$
\end{lemma}
\begin{proof}
Aus dem binomischen Lehrsatz(\ref{Th_3}) folgt, dass der Koeffizient von $X^i$ in dem Polynom ${n \choose i} a^{n-i}$ ist.\newline\newline
$"\Rightarrow"$\newline
Angenommen n ist eine Primzahl, dann ist es nach (\ref{th_25}) klar, dass $\forall i $, 0 < i < n,\newline\smallskip ${n \choose i} = \frac{n!}{(n-i)! i!} = 0 $ (mod n). Das heißt alle Koeffizienten sind Null.\newline\smallskip Für i = 0 erhält man  ${n \choose 0} a^n X^0 = a^n$, analog für $i = n$, ${n \choose n} a^0 X^n = X^n$. Daraus folgt:
$(X + a)^n = a^n + 0 + 0 + ... + 0 + X^n = a^n + X^n$(mod n).\newline\newline
$"\Leftarrow"$\newline
Sei $n$ nun eine zusammengesetzte Zahl(COMPOSITE). Betrachte einen Faktor q von n, mit der Vielfachheit k(Dabei ist zu beachten, dass Für 1 < q < n, $q^k | n$, aber $q^{k+1} \nmid n$).\newline
Der Koeffizient von $X^q$ sieht wie folgt aus:\newline\smallskip
\begin{equation}
    {n \choose q} \cdot a^{n-q} = \frac{n!}{(n-q)! n!} \cdot a^{n-q} = \frac{n(n-1)\cdot \cdot \cdot (n-q+1)}{q!} \cdot a^{n-q}.
\end{equation}
\newline\newline
Im Nenner lässt sich $q!$ als $q \cdot (q-1)!$ schreiben und im Zähler lässt sich n als $q^k\cdot m$ schreiben, $m \in \mathbb{Z}_{+}$. Das q im Nenner hebt sich mit einem der qs im Zähler auf. Der resultierende Term ist daher nicht durch $q^k$ teilbar, außerdem sind $q^k$ und $a^{n-k}$ teilerfremd. Daraus folgt $(X + a)^n \neq X^n + a $(mod $n$).
\end{proof}

Es wäre nun möglich anhand dieser Identität einen Primzahltest für eine Zahl $n$ zu realisieren. Dies wäre jedoch sehr ineffizient, da im Polynom die Auswertung von $n$ Koeffizienten nötig ist. Mit anderen Worten der Algorithmus braucht $\Omega(n)$ um zu entscheiden, ob die Zahl $n$ eine Primzahl ist oder nicht, und das ist nicht in Polynomialzeit realisierbar.\newline\newline Die Idee von AKS ist nicht nur modulo $n$, sondern auch modulo ein Polynom ($X^r -1$) zu nehmen. Das bedeutet, dass wir nicht nur alle Koeffizienten $c$ durch $c$ modulo $n$ ersetzen, sondern auch jedes $X^k$ durch $X^k mod \, r$ ersetzen, also wir arbeiten immer mit Polynomen vom Grad weniger als $r$. Dabei wird $r$ so gewählt, dass die Anzahl an Berechnungen kleiner als bei (\ref{eq:1}) ist, da der Rest von $(X + a)^n$(mod $ n, X^r - 1$) nur $ r + 1$ Koeffizienten hat.\newline Dies kann für ein entsprechend kleines $r$ in Polynomialzeit berechnet werden. Daher das Hauptziel jetzt ist ein entsprechend kleines $r$ zu wählen und zu testen, ob die Gleichung:\newline\newline
\begin{equation}\label{eq:2}
    \centerline{$(X + a)^n = X^n + a $(mod $X^r - 1, n$)}
\end{equation}

erfüllt ist.\newline

Nach Lemma \ref{hauptlemma} ist die Gleichung (\ref{eq:2}) für alle Primzahlen erfüllt. Aber ein Problem bei diesem Ansatz wäre, dass es auch zusammengesetzte Zahlen gibt, für die die Gleichung für manche Werte von $a$ und $r$ erfüllt ist. Jedoch ist das geeignete $r$ von oben durch ein Polynom in $log n$ beschränkt, das heißt der Algorithmus muss nur $log n$ a testen, um eine Entscheidung über die Primalität der Zahl $n$ zu treffen.

\end{flushleft}
\subsection{Der Algorithmus}\label{algo}
\textbf{Notation:}\newline
\textbf{Schritt 2: }$o_r(n)$ bezeichnet die Ordnung der
Zahl(Def. \ref{Df_4}).\newline \textbf{Schritt 3:}(a,b) ist der größte gemeinsame Teiler (Def. \ref{Df_1})
\newline \textbf{Schritt 5:}$\phi(n)$ ist die eulersche Phi-Funktion(Def. \ref{Df_6})  
%to do: schreibe den AKS Algorithmus
\begin{algorithm}[H]
\SetAlgoLined
\KwIn{$n \in \mathbb{N}, n \geq 2$.}

\begin{enumerate}
% STEP 1
\item \textbf{if} $n = a^b, a \in \mathbb{N}, b \geq 1$ , \textbf{return} COMPOSITE.
%STEP 2
\item  finde das kleinste $r$, sodass $o_{r}(n) > log^2 n $.
% STEP 3
\item \textbf{if} $1 < (a,n) < n, a \geq n $, \textbf{return} COMPOSITE.
%STEP 4
\item \textbf{if} $n \leq r $, \textbf{return} PRIME.
%STEP 5
\item \textbf{for} a = 1 to $\lfloor \sqrt{\phi(r)}log(n) \rfloor$:

 \textbf{if}$(X + a)^n \neq X^n + a $(mod $X^r - 1, n$), \textbf{return} COMPOSITE.
 %STEP 6
 \item \textbf{return} PRIME.
\end{enumerate}
 
\caption{AKS-Primzahltest}
\end{algorithm}

Agrawal, Kayal und Saxena haben drei Kriterien gefunden, wenn sie für eine Zahl erfüllt sind, dann ist diese Zahl eine Primzahl. 

\begin{prop}\label{aks-prop}
Seien $n,r \geq 2$  natürliche Zahlen, mit $ r \leq n$ und $o_{r}(n) > log^2 n$. $n$ ist genau dann prim, wenn:
\begin{enumerate}
    \item $n \neq a^b, \, für \,  b,a \in \mathbb{N}$.
    \item $n$ keinen Primfaktor $p \leq r$ hat.
    \item $(X + a)^n = X^n + a (mod \, n, \, X^r - 1)$ für alle $ a, \, 1 \leq a \leq \sqrt{r} log n$
\end{enumerate}
\end{prop}

Der Algorithmus überprüft alle drei Bedingungen für eine Zahl $n$.\newline

\textbf{Schritt 1: }im ersten Schritt sucht der Algorithmus zwei Zahlen $a,b \in \mathbb{N}, \, b > 1$, sodass $n = a^b$. Wenn solche Zahlen existieren, gibt er COMPOSITE zurück.  

\textbf{Schritt 2: }Hier probiert der Algorithmus sukzessive Werte von $r$ aus, bis $n^k \neq 1 (mod \, r), \, \forall k \leq log^2 n$.

\textbf{Schritt 3: } in diesem Schritt wird überprüft, ob $n$ gemeinsame Faktoren mit $a \leq r$ hat. 

\textbf{Schritt 4: } wenn keine Primfaktoren bei Schritt 3 gefunden worden sind, gibt der Algorithmus bei diesem Schritt PRIME zurück. Da $n \leq r$ und $(k,n) = 1, \, \forall k \leq r $(sonst hätte er beim 3. Schritt schon COMPOSITE zurückgegeben).

\textbf{Schritt 5: } hier wird die dritte Bedingung aus \ref{aks-prop} überprüft.

\textbf{Schritt 6: } wenn kein $a$ in Schritt 5 gefunden wurde, dann muss die Zahl prim sein. 
\subsection{Korrektheitsbeweis}

\begin{theorem}
Der Algorithmus gibt genau dann PRIME zurück, wenn $n$ eine Primzahl ist.
\end{theorem}
Der Korrektheitsbeweis lässt sich in zwei Teilen zerlegen, der erste Teil befasst sich mit der Hinrichtung des Beweises. Also, dass der Algorithmus PRIME liefert, wenn die Eingabe eine Primzahl ist. Dies ist aber trivial und braucht keine Lemmas. Für den zweiten Teil(Rückrichtung) sind jedoch mehrere Lemmas und Sätze nötig. Diese werden in diesem Kapitel in der Reihenfolge, wie sie im Orginalpaper vorkamen, dargestellt und bewiesen. Zudem wird vor jedem Lemma kurz erläutert, warum dieses Lemma für den Algorithmus von Bedeutung ist. Das erste Lemma ist eine Eigenschaft für das kleinste gemeinsame Vielfaches von $m$ Zahlen.

Das zweite Lemma beschäftigt sich mit der Existenz und der oberen Schranke vom $r$, das in der Grundidee erwähnt wurde.

Danach wird eine kurze Einführung zu introspektiven Zahlen und Polynomen dargestellt, dort wird der Begriff introspektiv definiert und genauer erklärt. Darüber hinaus werden Eigenschaften introspektiver Polynome und Zahlen bewiesen.

Schließlich werden zwei sehr wichtige Mengen definiert.% do some more stuff here

\textbf{Erster Teil des Beweises($\Rightarrow$)}
\begin{theorem}
Wenn $n$ eine Primzahl ist, dann gibt der Algorithmus PRIME zurück.
\end{theorem}

\begin{proof}
Wenn $n$ eine Primzahl ist, gibt der erste Schritt in (\ref{algo}) niemals COMPOSITE zurück, da keine Primzahl sich als $a^b$ schreiben lässt, sonst wäre $n$ durch $a$ teilbar. Der dritte Schritt kann auch niemals COMPOSITE zurückgeben, da $(a,n) = 1, \forall a$. Nach Lemma (\ref{hauptlemma}) kann Schritt 5 auch nie COMPOSITE zurückgeben. Daher muss der Algorithmus entweder bei viertem oder bei sechstem Schritt PRIME zurückgeben.  
\end{proof}

\textbf{Zweiter Teil des Beweises($\Leftarrow$)}

Hier wird die Rückrichtung bewiesen; also  wenn die Ausgabe PRIME ist, dann ist $n$ eine Primzahl.
Wenn der Algorithmus bei Schritt 4 PRIME zurückgibt, dann muss $n$ eine Primzahl sein, da sonst der Algorithmus einen nicht trivialen Faktor von $n$ in Schritt 3 gefunden hätte. Das heißt, es bleibt nur der Fall bei Schritt 6. Der Algorithmus hat Zwei Hauptschritte, das sind Schritt 2 und Schritt 5. Bei Schritt 2 wird ein geeignetes $r$ gefunden und bei Schritt 5 wird verifiziert, ob die Gleichung (\ref{eq:2}) für mehrere a Werte gilt.    

Für diesen Teil ist eine Vorarbeit nötig. Das erste wichtige Lemma ist die Abschätzung des kleinsten gemeinsamen Vielfaches.


\begin{lemma}\label{lemma_1}
Für $m \geq 7$:\newline
\begin{equation}
    LCM(m) \geq 2^m.  
\end{equation}
\end{lemma}

\begin{proof}
Sei $d_{n} = LCM_{1 \leq m \leq n}{\{m\}}$, betrachte das folgende Integral für $n \geq 1$:\newline
\begin{equation}\label{intg}
    \begin{split}
      \begin{aligned}
        I_{n,m}&= \int_{0}^{1} x^{m-1} (1-x)^{n - m} dx \underbrace{=}_{(\ref{Th_3})} \int_{0}^{1} x^{m-1} \sum_{r = 0}^{n - m} (-1)^r {n - m\choose r} \cdot x^r dx \\
        &= \sum_{r = 0}^{n - m}(-1)^r {n - m \choose r} \cdot \int_{0}^{1} x^{m + r - 1} 
        = \sum_{r = 0}^{n - m} (-1)^r {n - m \choose r} \cdot \frac{x^{m+r}}{m+r}\Big|_0^1\\
        &= \sum_{r = 0}^{n - m} (-1)^r {n - m \choose r} \cdot \frac{1}{m + r}
      \end{aligned}
    \end{split}
\end{equation}

Aus (\ref{intg}) ist es leicht zu sehen, $r \leq n - m $ und folglich $ r + m \leq n $. Das heißt $m + r \mid d_{n}$. Dabei ist es offensichtlich, dass $d_{n} \cdot I_{n,m} \in \mathbb{N}$, für $1 \leq m \leq n$.

Durch Induktion nach $n - m$ und partielle Integration lässt sich zeigen, dass für $n \in \mathbb{N}$ und m, mit $1 \leq m \leq n$:
\begin{equation}\label{I_mn}
    I_{n,m} = \frac{1}{m \cdot {n \choose m}}
\end{equation}

gilt.\newline\newline


Da $m + r \mid d_{n} $ und (\ref{I_mn}) gelten, muss
\begin{equation}
    m \cdot {n \choose m} \mid d_{n}
\end{equation}
$\forall m $ mit $1 \leq m \leq n $ auch gelten.\newline

Somit gilt auch: 
\begin{equation}\label{d2n}
    n {2n\choose n} \mid d_{2n}
\end{equation}

Beziehungsweise 

\begin{equation}\label{d2n1}
    (2n + 1) {2n \choose n} = (n + 1) {2n + 1 \choose n + 1} \mid d_{2n + 1} 
\end{equation}

\smallskip
Aus (\ref{d2n}) und (\ref{d2n1}) folgt $d_{n} \mid d_{n+1}$\newline\newline Sowie 
\begin{equation}\label{helps}
    n(2n + 1) {2n \choose n} \mid d_{2n + 1}
\end{equation}
\newline\newline
Es bleibt zu zeigen, dass $d_{2n+1} $eine obere Schranke von $2^{2n + 1}$ ist.\newline\newline  
$\Rightarrow d_{2n + 1} \underbrace{\geq}_{(\ref{helps})} n (2n + 1) {2n \choose n} \geq n \cdot 4^n \geq 2 \cdot 2^{2n} = 2^{2n + 2}$ für $n \geq 2$.\newline\newline
Für $n \geq 4$ ist $d_{2n + 2} \geq d_{2n + 1} \geq 2^{2n + n }$, wenn diese Aussage für $n \geq 4 $ gilt, dann gilt sie auch für $n \geq 7$.\newline\newline Generell gilt: $d_{m} = LCM_{1 \leq m \leq n} \{m\} = LCM(m) \geq  2^m $.
\end{proof}

\textbf{Abschätzung von r}.\newline\newline
Das zweite Lemma befasst sich mit der Abschätzung von $r$, das vorher bei der Grundidee erwähnt wurde. Dabei wird gezeigt, dass $r$ eine logarithmische obere Schranke hat, und das ermöglicht die Polynomiallaufzeit.  

\begin{lemma}
Es existiert ein $ r \leq max \{ 3, \lceil log^5 n \rceil \}$, sodass $o_{r}(n) > log^2 n$.
\end{lemma}

\begin{proof}
Der Beweis lässt sich in 3 Schritten zerlegen; zeigen, dass ein $r \leq B$ existiert, danach wird die Existenz von $o_{r}(n)$ gezeigt, als letztes wird die Eigenschaft $o_{r}(n) \geq log^2 n$ für dieses $r$ bewiesen. 
Sei $n \geq 1$.\newline\newline Für $n = 2$ und $r = 3$ gilt trivialerweise :\newline\newline $ord_{3}(2) = 2^2 = 4 = 1 $(mod 3) > 1 = $log^2 2$.

Ab jetzt wird angenommen, dass $n > 2$.\newline\newline
\textbf{Bemerkung: } $\lceil log^5 3 \rceil = 11 \Rightarrow log^5 n > 10, \forall n > 2$.\newline\newline

Sei $B = \lceil log^5 n \rceil$, nach Lemma \ref{lemma_1} gilt $LCM(B) \geq 2^B$. Als erstes ist zu zeigen: es existiert ein $r \leq B$, sodass
\begin{equation}\label{prod_cor}
     N_{B} = n^{\lfloor log B \rfloor } \cdot \prod_{i = 1}^{\lfloor log^2 n \rfloor} (n^i - 1)
\end{equation}
$r$ dieses Produkt nicht teilt. Dies kann durch Widerspruch bewiesen werden.

Angenommen, $\forall 1 \leq r \leq n$, $r$ teilt $N_{B}$. Das Produkt ist ein gemeinsames Vielfaches aller Zahlen kleiner gleich $B$, daher muss $N_{B} \geq LCM(B)$ gelten. Jedoch hat $N_{B}$ eine obere Schranke, betrachte das folgende:
\begin{equation}\label{long_p}
    \begin{split}
        \begin{aligned}
            N_{B} &= n^{\lfloor log B \rfloor } \cdot \prod_{i = 1}^{\lfloor log^2 n \rfloor} (n^i - 1)
                < n^{\lfloor log B \rfloor } \cdot \prod_{i = 1}^{\lfloor log^2 n \rfloor} n^i \\
                &= n^{\lfloor log B \rfloor + \sum_{i = 1}^{\lfloor log^2 n \rfloor } i}\\ 
                &\underbrace{=}_{(\ref{gauss})} n^{\lfloor log B \rfloor + \frac{log^2 n \cdot (log^2 n + 1)}{2}}  \\
                &= n^{\lfloor log B \rfloor + \frac{log^4 n + log^2 n}{2}}\\
                & \leq n^{\lfloor log B \rfloor + \frac{log^4 n + \frac{log^4 n}{2}}{2}}\\
                & \leq n^{\lfloor \frac{log^4 n}{4} \rfloor + \frac{log^4 n + \frac{log^4 n}{2}}{2}}\\
                & \leq n^{log^4 n} = 2^{log n^{log^4 n}} = 2^{log^5 n} = 2^{B}.
        \end{aligned}
    \end{split}
\end{equation}
Aus (\ref{long_p}) folgt: $LCM(B) \leq N_{B} \leq 2^B$. Das ist aber ein Widerspruch zu Lemma \ref{lemma_1}. Das heißt es existiert eine Menge von Zahlen $R = \{r_{1},r_{2},...,r_{t}\}$, $ 1 \leq r_i \leq B$,$\, i = 1,2, ..,t$, sodass $r_{i}$ das Produkt nicht teilt. Ferner sei $r_{m}$ das kleinste Element dieser Menge.\newline\newline Jetzt bleibt zu zeigen, dass $o_{r_{m}}(n)$ für dieses $r_{m}$ existiert, und $o_{r_{m}}(n) \geq log^2 n$.\newline\newline Sei $r_{m} = ab $, wobei $a$ aus den Primfaktoren besteht, die $n$ teilen und $b$ aus den restlichen Primfaktoren. Offensichtlich $(b,n) = 1$, außerdem ist die höchste Potenz, die ein Primfaktor(bei der Primfaktorzerlegung) haben kann, kleiner als $\lfloor log B \rfloor$. Da sonst $a$ größer als $B$ wäre. Das heißt, jeder Primfaktor in $a$ hat einen Exponent kleiner als der Exponent beim gleichen Primfaktor von $n$ und jeder vorkommende Faktor in $a$, kommt auch in $n$ vor(alle Primfaktoren von $a$ teilen $n$).\newline
$\Rightarrow a \mid n^{log B}$.\newline
Das heißt $b$ teilt $\prod_{i = 1}^{\lfloor log^2 n \rfloor}(n^i - 1)$ nicht, da sonst $r_{m}$ das Produkt teilen würde. $b$ teilt $n^{log B}$ auch nicht, da $n$ und $b$ keine gemeinsamen Primfaktoren haben($(b,n) = 1$). Zudem war $r_{m}$ das kleinste Element aus $R$, das $N_{B}$ teilt, daher ist $r_{m}= b$ und $(b,n) = 1$ beziehungsweise $(r_{m},n) = 1$. Somit existiert $o_{r_{m}}(n)$.\newline\newline

Zuletzt bleibt nur der Beweis der Ungleichung $o_{r_{m}}(n) > log^2 n$. Das lässt sich auch wie oben durch einen Widerspruch zeigen. Sei $o_{r_{m}}(n) = k \leq log^2 n$, per Definition ist $k$ die kleinste Zahl, für die $n^k = 1$ (mod $r_{m}$) wahr ist. daher gilt auch:\newline\newline
\centerline{$n^k - 1 = 0 $(mod  $r_{m}$) $\Rightarrow r_{m} \mid n^k - 1$.}

Aber wenn das gelten würde, würde $r_{m}$ $N_{B}$ teilen (z.B. $ r_{m} \mid n^{log^2 n} -1$) und das führt zum Widerspruch, da $r_{m} \nmid N_{B}$. 
\end{proof}

Bevor wir mit dem Beweis fortfahren, brauchen wir das folgende Lemma:

\begin{lemma}\label{ord_prime_l}
wenn $o_{r}(n) > 1$, dann existiert ein Primfaktor $p$ von $n$, sodass $o_{r}(p) > 1$ gilt. 
\end{lemma}
\begin{proof}
Sei $o_{r}(n) > 1$, mit der Primfaktorzerlegung $n = \prod_{i = 1}^{M} p_{i}^{e_{i}}$.\newline\newline Weiterhin sei $o_{r}(p_{i}) = 1$,$\forall 1 \leq i \leq M$ (es existiert kein $p_{i} : o_{r}(p_{i}) > 1$).\newline\newline
$\Rightarrow p_{i} = 1 $(mod $r$), $\forall i$. Folglich ist $n = \prod_{i = 1} ^{M} p_{i}^{e_{i}} = \prod_{i = 1} ^{M}(1)^{e_{i}} = 1 $(mod $r$).\newline\newline Daraus wird ersichtlich, dass $o_{r}(n) = 1$ und das ist ein Widerspruch zur Voraussetzung, dass $o_{r}(n) > 1$. Daher muss mindestens ein Primfaktor $p : o_{r}(p) > 1$.  
\end{proof}

\smallskip

Nach Lemma (\ref{ord_prime_l}) muss für  $o_{r}(n) > 1$ ein Primteiler $p$ von $n$(Aus der Primfaktorzerlegung) existieren, sodass $o_{r}(p) > 1$. Außerdem ist es klar, dass $p > r$, da sonst der Algorithmus eine Entscheidung über die Primalität schon im 3. beziehungsweise im 4. Schritt getroffen hätte(also, für den Fall $a \leq r$ oder $n \leq r$). Seien $p,n \in \mathbb{Z}_{r}^{*}$ fest. Ferner sei  $l = \lfloor \sqrt{\phi(r)} log n \rfloor$.\newline\newline

Der 5. Schritt des Algorithmus überprüft $l$ Gleichungen. Wir wollen nun die Ausgabe PRIME untersuchen, also nach Voraussetzung ist die Ausgabe PRIME . Das heißt, für alle Gleichungen $l$ gilt:\newline\newline
\centerline{$(X + a)^n = X^n + a $ (mod $X^r - 1, n$)}.

$\forall a$, $ 0 \leq a \leq l$.\newline\newline
\textbf{Bemerkung:} Für a = 0 $\Rightarrow (X + 0)^n = X^n + 0$ und damit ist die Gleichung trivialerweise wahr.\newline\newline 
Aus der oberen Gleichung kann man drei andere Gleichungen folgern:
\begin{enumerate}
    \item $(X + a)^n = X^n + a $(mod $X^r - 1, p$).
$\forall$ $ 0 \leq a \leq l$.\newline\newline  
Da $p$ ein Primfaktor von $n$, gilt auch nach Lemma \ref{hauptlemma} \newline
\item $(X + a)^p = X^p + a $(mod $X^r - 1, p$).
$\forall$ $0 \leq a \leq l$.\newline\newline
Aus 1 und 2 gilt auch\newline
\item $(X + a)^{\frac{n}{p}} = X^{\frac{n}{p}} + a $(mod $X^r - 1, p$).
$\forall$ $0 \leq a \leq l$.\newline\newline
\end{enumerate}

Das heißt,$\forall$ $0 \leq a \leq l$, $n$, $\frac{n}{p}$ und $p$ haben ein identisches Verhalten bezüglich der obigen Gleichung. Diese Eigenschaft heißt Introspektivität.\newline    


\textbf{Introspektive Zahlen und Polynome:}
Introspektive Zahlen und Polynome spielen eine große Rolle für den Korrektheitsbeweis des AKS-Primzahltests. Zunächst ist eine formale Definition von Introspektiven Zahlen beziehungsweise Polynomen nötig.

\begin{definition}
Seien $f(X)$ ein Polynom in $\mathbb{Z}[X]$ und $m \in \mathbb{N}$ , $m$ ist \textbf{introspektiv} bezüglich $f(X)$, wenn\newline\newline
    \centerline{$[f(X)]^{m} = f(X^m)$ (mod $X^r - 1,p$)} 
gilt. 
\end{definition}

Für den Korrektheitsbeweis ist eine Eigenschaft introspektiver Zahlen beziehungsweise Polynome relevant. Dies ist die Multiplikativität von solchen Zahlen und Polynomen.\newline\newline

\begin{lemma}\label{intros_num}
Seien $m,m^{'}$ introspektive Zahlen bezüglich eines Polynoms $f(X)$ in $\mathbb{Z}[X]$. Dann ist $m \cdot m^{'}$ auch introspektiv bezüglich $f(X)$.
\end{lemma}

\begin{proof}
Da $m$ bezüglich $f(X)$ introspektiv ist, gilt:\newline\newline
\centerline{$[f(X)]^{m \cdot m^{'}} = [f(X)^{m}]^{m^{'}}$( mod $X^r - 1, p$).}\newline\newline
$m^{'}$ ist auch introspektiv bezüglich $f(X^{m})$, das heißt es gilt auch:\newline\newline
\centerline{$[f(X^m)]^{m^{'}} = f(X^{m \cdot m^{'}})$( mod $X^{m\cdot r} - 1, p$).}\newline\newline

Offensichtlich gilt $X^{m \cdot r } - 1 = (X^{r})^{m} - (1)^m \underbrace{=}_{(\ref{funny_id})} (X^{r} - 1) \sum_{k = 0}^{m-1} X^{r \cdot k}$.\newline\newline$\Rightarrow X^{r} - 1\mid X^{mr} - 1$. Somit gilt auch\newline\newline
\centerline{$[f(X^m)]^{m^{'}} = f(X^{m \cdot m^{'}})$( mod $X^{r} - 1,p$).}\newline

Daraus folgt:\newline\newline
\centerline{$[f(X)]^{m \cdot m^{'}} = f(X^{m \cdot m^{'}}) $( mod $X^r - 1, p$).}\newline

Das heißt wenn eine Zahl $m$ bezüglich eines Polynoms $f(X)$ introspektiv ist, ist auch $m \cdot m^{'}$ bezüglich des selben Polynoms $f(X)$ introspektiv. 
\end{proof}

\smallskip

\begin{lemma}
wenn $m$ bezüglich der Polynome $f(X)$ und $g(X)$ introspektiv ist. Dann ist $m$ auch bezüglich $f(X) \cdot g(X)$ introspektiv.  
\end{lemma}


\begin{proof}\label{intros_pol}
    \centerline{$[f(X) \cdot g(X)]^{m} = [f(X)]^m \cdot [g(X)]^m \underbrace{=}_{(\ref{intros_num})} f(X^m) \cdot g(X^m)$(mod $X^r - 1, p$).}
\end{proof}

Für den Beweis benötigen wir zwei Gruppen, wir definieren nun die folgenden Mengen:\newline\newline
\begin{equation}\label{I_Set}
      I := \{ \frac{n^i}{p^i} \cdot p^j \mid i,j \geq 0\}.
\end{equation}
\begin{equation}
     P := \{ \prod_{a = 0}^{l} (X + a)^{e_{a}} \mid e_{a} \geq 0 \}.
\end{equation}\label{P_Set}
\newline\newline
Aus Lemma \ref{intros_num} und Lemma \ref{intros_pol} folgt, dass jedes Element $i \in I$ bezüglich jedes Polynoms $p \in P$ introspektiv ist.\newline\newline

Wir definieren jetzt zwei Gruppen basierend auf diese zwei Mengen, die für den Beweis von Bedeutung sind.\newline\newline

Die erste Gruppe G sei die Gruppe aller Restklassen in $I$ modulo $r$. Da $(n,r) = (p,r) = 1 $ ist G eine multiplikative Untergruppe von $\mathbb{Z}_{r}^{*}$. Sei die Kardinalität im weiteren $ | G | = t $. G wird von $n$ und $p$ modulo $r$ erzeugt, und da $o_{r}(n) > log^2 n$, gilt auch $t > log^2 n$.\newline\newline

Um die zweite Gruppe zu definieren, sind ein paar Eigenschaften von zyklotomischen Polynomen über endlichen Körpern erforderlich.

\begin{lemma}[\textbf{Irreduzible Teiler zyklotomischer Polynome}] Sei $\mathbb{K}_{p}$ ein Körper und $Q_{r}(X)$ das r-te zyklotomische Polynom über $\mathbb{K}_{p}$. Dann existiert ein irreduzibler Teiler $h(X)$ von $X^r - 1$ über $\mathbb{K}_{p}$ vom Grad $o_{r}(p).$
\end{lemma}

\begin{proof}
    
\end{proof}

\section{Laufzeitanalyse}

\section{Zusammenfassung}


%%%%%%%%%%%%%%%%%%%%%%%%%%%%
%% Literaturverzeichnis wird 
%% automatisch eingefügt
%%%%%%%%%%%%%%%%%%%%%%%%%%%%
\clearpage
\lhead{}
\printbibliography
\addcontentsline{toc}{section}{\bibname}


%%%%%%%%%%%%%%%%%%%%%%%%%%%%
%% Anhang (optional) 
%%%%%%%%%%%%%%%%%%%%%%%%%%%%
\clearpage
\appendix
\section{Anhang A}

%%%%%%%%%%%%%%%%%%%%%%%%%%%%
%% Eidesstattliche Erklärung
%% muss angepasst werden 
%% in Erklaerung.tex
%%%%%%%%%%%%%%%%%%%%%%%%%%%%
\newpage
\begin{otherlanguage}{ngerman}
\thispagestyle{empty}
\section*{Eidesstattliche Erklärung}
\thispagestyle{empty}
Hiermit versichere ich, die vorliegende Arbeit selbstständig verfasst und keine anderen als die angegebenen Quellen und Hilfsmittel benutzt sowie die Zitate deutlich kenntlich gemacht zu haben.
\newline
\vspace{4\baselineskip}\\
Leipzig, den \today \hfill Salman Salman 
\vspace{4\baselineskip}\\
\end{otherlanguage}

\end{document}
